\documentclass[a4paper,11pt]{article}

\usepackage{amsmath}		
\usepackage{amssymb}		
\usepackage{siunitx}		
\usepackage{graphicx}
\usepackage{geometry}

\usepackage{aas_macros}
\usepackage[
backend=biber,
style=ieee,
sorting=none
]{biblatex}
\addbibresource{References.bib}
\setlength{\biblabelsep}{.5em}

\usepackage{hyperref}			
\hypersetup{
    colorlinks=true,
    linkcolor=blue,
    filecolor=magenta,      
    urlcolor=blue,
    citecolor=cyan,
}
\newgeometry{
    top=0.75in,
    bottom=0.75in,
    outer=0.75in,
    inner=0.75in,
}
\usepackage{lipsum}
\begin{document}

\twocolumn[
\title{}
\begin{abstract}
Overall scientific questions.
\end{abstract}
\vspace{0.5cm}
]

% -----------------------------------------------

\section*{Introduction}

\begin{itemize}
    \item Overview of the ISM.\ How dust forms, why we use CO as a tracer. Why we might use other tracers.
    \item Jeans collapse a model for star formation. 
    \item IMF / CMF.\ How the high mass end changes. SF across different environments.
    \item Quasistatic turbulent core fragmentation as a model for star formation.
    \item Differences in methods.
    \item Literature on decoupling cores from clumps. We want to know at what stage the gas decouples.
    \item CHIMPS paper and other extragalactic attempts at this sort of work.
\end{itemize}

\par\noindent The interstellar medium (ISM) is made up of primarily gas, as well as dust and cosmic rays. It consists of 5 main phases, one of which will be discussed here. The cold neutral medium 

% -----------------------------------------------

\section*{Project Work}

\begin{itemize}
    \item Motivations (link back to the lit review).
    \item Methods.
    \item Conclusions and what that might mean.
    \item Future work. 
\end{itemize}

\par\noindent

% -----------------------------------------------

\par\noindent\cite{SEDIGISM}

\section*{Bibliography}
% \hfuzz=1.5pt
\printbibliography[heading=none]

\end{document}